% !TeX root= ../quantum.tex


\chapter*{پیشگفتار}
\label{chap:preface}

اطلاعات کوانتومی\LTRfootnote{Quantum information} به استفاده از مفاهیم مکانیک کوانتومی برای انجام پردازش اطلاعات\LTRfootnote{Information processing} و انتقال اطلاعات\LTRfootnote{Information transmission} مربوط می‌شود. پردازش اطلاعات کوانتومی\LTRfootnote{Quantum information processing (QIP)} (\lr{QIP}) یک حوزه پژوهشی هیجان‌انگیز با کاربردهای متعدد است که از جمله آن‌ها می‌توان به توزیع کلید کوانتومی\LTRfootnote{Quantum key distribution (QKD)} (\lr{QKD})، دورنوردی کوانتومی\LTRfootnote{Quantum teleportation}، محاسبات کوانتومی و محاسبات کوانتومی توزیع‌شده\LTRfootnote{Distributed quantum computing}، حسگری کوانتومی\LTRfootnote{Quantum sensing}، شبکه‌بندی کوانتومی\LTRfootnote{Quantum networking}، رادار کوانتومی\LTRfootnote{Quantum radar} و یادگیری ماشین کوانتومی\LTRfootnote{Quantum machine learning} اشاره کرد. علم اطلاعات کوانتومی\LTRfootnote{Quantum information science (QIS)} (\lr{QIS})، با توجه به تعداد کتاب‌های منتشر شده در زمینه‌های \lr{QIP}، محاسبات کوانتومی و ارتباطات کوانتومی و همچنین کنفرانس‌های متعددی که صرفاً به \lr{QIS} اختصاص یافته‌اند، در حال تجربه توسعه سریعی است. با توجه به نوظهور بودن مفاهیم زیربنایی \lr{QIS}، انتظار می‌رود این موضوع مورد توجه بسیاری از دانشمندان، مهندسان، آژانس‌های دولتی، صنایع و سرمایه‌گذاران قرار گیرد، و نه فقط کسانی که در پژوهش‌های \lr{QIP} درگیر هستند. درهم‌تنیدگی\LTRfootnote{Entanglement} نشان‌دهنده یک ویژگی منحصر‌به‌فرد در \lr{QIP} است که کامپیوترهای کوانتومی را قادر می‌سازد مسائلی را حل کنند که از نظر عددی برای کامپیوترهای کلاسیک دشوار\LTRfootnote{Numerically intractable} هستند. این ویژگی همچنین حسگرهای بهبودیافته کوانتومی\LTRfootnote{Quantum-enhanced sensors} را با حساسیت‌های اندازه‌گیری\LTRfootnote{Measurement sensitivities} فراتر از محدودیت‌های کلاسیک توانمند ساخته و امنیت انتقال داده‌ی تصدیق‌پذیر\LTRfootnote{Certifiable data transmission security} را فراهم می‌کند که توسط قوانین مکانیک کوانتومی تضمین می‌شود؛ برخلاف فرضیات اثبات‌نشده‌ای که در رمزنگاری متداول\LTRfootnote{Conventional cryptography} برای امنیت محاسباتی\LTRfootnote{Computational security} استفاده می‌شوند. از درهم‌تنیدگی از پیش اشتراک‌گذاری شده\LTRfootnote{Preshared entanglement} می‌توان برای موارد زیر استفاده کرد: (۱) بهبود ظرفیت کانال کلاسیک\LTRfootnote{Classical channel capacity}، (۲) فراهم‌سازی ارتباطات امن، (۳) بهبود حساسیت حسگر، (۴) توانمندسازی محاسبات کوانتومی توزیع‌شده، (۵) فراهم‌سازی حسگری توزیع‌شده به کمک درهم‌تنیدگی\LTRfootnote{Entanglement-assisted distributed sensing}، و (۶) فراهم‌سازی دسترسی به کامپیوتر کوانتومی با امنیت اثبات‌پذیر، و همچنین انجام بسیاری از وظایف دیگر.

به دلیل ماهیت بین‌رشته‌ای \lr{QIS}، این کتاب با هدف ایجاد تعادل مناسب میان \lr{QIP}، تصحیح خطای کوانتومی\LTRfootnote{Quantum error correction}، ارتباطات کوانتومی، حسگری کوانتومی و شبکه‌بندی کوانتومی نگاشته شده است. اهداف اصلی کتاب را می‌توان به شرح زیر خلاصه کرد:

\begin{enumerate}
	\item کتاب روندها را در \lr{QIP}، تصحیح خطای کوانتومی، ارتباطات کوانتومی، حسگری کوانتومی، شبکه‌بندی کوانتومی و یادگیری ماشین کوانتومی توصیف می‌کند.
	\item کتاب یک مقدمه جامع و خودکفا\LTRfootnote{Self-contained introduction} برای ارتباطات کوانتومی، تصحیح خطای کوانتومی، حسگری کوانتومی و شبکه‌بندی کوانتومی ارائه می‌دهد.
	\item کتاب طیف گسترده‌ای از خوانندگان را هدف قرار می‌دهد: مهندسان برق، مهندسان مخابرات، مهندسان اپتیک، ریاضی‌دانان کاربردی، دانشمندان علوم کامپیوتر و فیزیک‌دانان.
	\item کتاب نیازی به دانش قبلی از مکانیک کوانتومی یا \lr{QIP} ندارد. مفاهیم پایه در فصل \ref{chap:3} ارائه شده‌اند.
	\item کتاب به هیچ پیش‌زمینه پیش‌نیازی نیاز ندارد، به‌جز درک مفاهیم پایه جبر برداری در سطح مقطع کارشناسی.
	\item برای خوانندگانی که با نظریه اطلاعات\LTRfootnote{Information theory}، کدگذاری کانال\LTRfootnote{Channel coding} و نظریه آشکارسازی\LTRfootnote{Detection theory} آشنا نیستند، فصل \ref{chap:2} مفاهیم و تعاریف پایه را در حوزه‌های مربوطه ارائه می‌دهد. این فصل تنها مفاهیم نظریه اطلاعات، کدگذاری و آشکارسازی را پوشش می‌دهد که در فصل‌های بعدی مرتبط هستند.
	\item خوانندگانی که به تصحیح خطای کوانتومی علاقه‌مند نیستند، همچنان می‌توانند از دیگر فصل‌های کتاب استفاده کنند.
	\item خوانندگانی که صرفاً به تصحیح خطای کوانتومی علاقه‌مند هستند، نیازی به مطالعه سایر فصل‌ها به‌جز فصل‌های \ref{chap:3} و \ref{chap:4} ندارند و می‌توانند تنها بر فصل‌های تصحیح خطای کوانتومی تمرکز کنند. در واقع، تلاشی صورت گرفته تا فصل‌ها مستقل باشند و در عین حال جریان مناسبی بین آن‌ها برقرار باشد.
	\item کتاب با مفاهیم شهودی\LTRfootnote{Intuitive concepts} و پایه شروع شده و سپس به سمت موضوعات با دشواری متوسط و در نهایت به موضوعات پیچیده و درگیر با ریاضیات حرکت می‌کند.
	\item چندین درس را می‌توان با استفاده از این کتاب ارائه داد. مثال‌ها شامل دروسی در زمینه‌های زیر است: (۱) تصحیح خطای کوانتومی، (۲) \lr{QIP} و یادگیری ماشین کوانتومی، (۳) ارتباطات کوانتومی، (۴) حسگری کوانتومی، (۵) شبکه‌بندی کوانتومی، و (۶) یکپارچه‌سازی\LTRfootnote{Integration} مفاهیم ارتباطات کوانتومی، حسگری کوانتومی و شبکه‌بندی کوانتومی.
	\item خواننده‌ای که این کتاب را به پایان برساند، قادر خواهد بود در زمینه‌های مختلف \lr{QIS} به پژوهش مستقل بپردازد.
\end{enumerate}

این کتاب یک مقدمه خودکفا برای اطلاعات کوانتومی، ارتباطات کوانتومی، تصحیح خطای کوانتومی، حسگری کوانتومی و شبکه‌بندی کوانتومی است. خواننده موفق کتاب برای مطالعه بیشتر در این حوزه‌ها آماده خواهد بود و برای انجام پژوهش مستقل مهیا می‌شود. خواننده‌ای که کتاب را به پایان رسانده باشد، قادر به طراحی مدارات \lr{QIP}، کدهای پایدارساز\LTRfootnote{Stabilizer codes}، کدهای \lr{CSS}\LTRfootnote{Calderbank-Shor-Steane (CSS) codes}، کدهای زیرسیستم\LTRfootnote{Subsystem codes}، کدهای توپولوژیکی\LTRfootnote{Topological codes}، کدهای سطحی\LTRfootnote{Surface codes} و کدهای تصحیح خطای کوانتومی به کمک درهم‌تنیدگی\LTRfootnote{Entanglement-assisted quantum error correction codes} خواهد بود. همچنین، خواننده در زمینه تصحیح خطای کوانتومی تحمل‌پذیر خطا\LTRfootnote{Fault-tolerant} مهارت خواهد یافت. علاوه بر این، خواننده در ارتباطات کوانتومی، \lr{QKD} و یادگیری ماشین کوانتومی مسلط خواهد شد و دانش لازم در حسگری کوانتومی و رادارهای کوانتومی را برای انجام پژوهش مستقل در این زمینه‌ها کسب خواهد کرد. خوانندگان این کتاب همچنین قادر خواهند بود فرستنده-گیرنده‌های کوانتومی\LTRfootnote{Quantum transceivers} مناسب برای کاربردهای ارتباطات کوانتومی و حسگری کوانتومی را پیشنهاد دهند. علاوه بر این، خوانندگان کتاب قادر به پیشنهاد پروتکل‌های شبکه‌بندی ارتباطات کوانتومی و معماری‌های شبکه‌بندی کوانتومی خواهند بود.

برخی از ویژگی‌های منحصر‌به‌فرد کتاب را می‌توان به شرح زیر خلاصه کرد:
\begin{enumerate}
	\item این کتاب حوزه‌های \lr{QIP}، ارتباطات، تصحیح خطا، حسگری، شبکه‌بندی و یادگیری ماشین را یکپارچه می‌کند.
	\item کتاب نیازی به دانش مکانیک کوانتومی ندارد.
	\item این کتاب به هیچ مطلب پیش‌نیازی به‌جز مفاهیم پایه جبر برداری در سطح کارشناسی نیاز ندارد.
	\item این کتاب شرح عمیقی از طراحی مدارها و ماژول‌های \lr{QIP}، تصحیح خطا، ارتباطات و شبکه‌بندی، و حسگری ارائه می‌دهد.
	\item دانشجوی موفق برای مطالعه بیشتر در این زمینه‌ها آماده شده و واجد شرایط انجام پژوهش مستقل خواهد بود.
	\item خواننده‌ای که این کتاب را به پایان رسانده باشد، قادر به طراحی مدارات پردازش اطلاعات، کدهای پایدارساز، کدهای \lr{CSS}، کدهای توپولوژیکی، کدهای سطحی، کدهای زیرسیستم، کدهای بررسی توازن کم‌چگال (\lr{LDPC}) کوانتومی\LTRfootnote{Quantum low-density parity-check (LDPC) codes} و کدهای تصحیح خطای کوانتومی به کمک درهم‌تنیدگی خواهد بود و در طراحی تحمل‌پذیر خطا مهارت خواهد یافت.
	\item دانشجوی موفق قادر خواهد بود فرستنده-گیرنده‌های کوانتومی را برای کاربردهای ارتباطات کوانتومی و حسگری کوانتومی پیشنهاد دهد.
	\item دانشجوی موفق قادر به پیشنهاد پروتکل‌های شبکه‌بندی ارتباطات کوانتومی و معماری‌های شبکه‌بندی کوانتومی خواهد بود.
	\item دانشجویانی که کتاب را به پایان رسانده باشند، در یادگیری ماشین کوانتومی مهارت خواهند یافت.
	\item در نهایت، دانشجویان در \lr{QKD} و دورنوردی کوانتومی مسلط خواهند شد.
	\item مطالب اضافی برای پشتیبانی از کتاب در یک وب‌سایت همراه که پس از انتشار توسعه می‌یابد، ارائه خواهد شد.
\end{enumerate}

نویسنده مایل است از حمایت بنیاد ملی علوم (\lr{NSF}) قدردانی نماید. در نهایت، تشکر ویژه‌ای از \lr{Prem Kumar Kaliamoorthi}، \lr{Kyle Gravel} و \lr{Tim Pitts} از انتشارات الزویر به‌خاطر تلاش‌های بی‌وقفه‌شان در سازماندهی امور لجستیکی کتاب، از جمله ویرایش و ترویج، که برای به ثمر رسیدن این کتاب ضروری بودند، به‌عمل می‌آید.









